% Part: sets-functions-relations
% Chapter: infinite
% Section: dedekind
%
\documentclass[../../../include/open-logic-section]{subfiles}

\begin{document}

\olfileid{sfr}{infinite}{dedekind}
\olsection{ڈیڈیکائنڈ الجبرے}

اسی قدرتی عدداں نوں صرف لامتناہی نہیں چاہندے؛ اسی چاہندے آں کہ اوہناں
کول کجھ (الجبری) خاصیتاں وی ہون: اوہناں نوں جمع، ضرب وغیرہ تحت درست
ورتاؤ کرنا چاہیدا اے۔

ڈیڈیکائنڈ دا خیال ایہ سی کہ \emph{جانشین فنکشن} دے خیال نوں بنیادی
منیا جاوے، تے فیر عدداں دی خاصیت ایہناں شرطاں نال دَسی جاوے:
\begin{enumerate}
	\item اک عدد $0$ اے جیہڑا کسے عدد دا جانشین نہیں
	\\i.e., $0 \notin \ran{s}$
	\\i.e., $\forall x\ s(x) \neq 0$
	\item وکھرے عدداں دے وکھرے جانشین ہوندے نیں
	\\i.e., $s$ اک !!a{injection} اے
	\\i.e., $\forall x \forall y (s(x) = s(y) \lif x = y)$
	\item\ollabel{repeatedapplication} ہر عدد، جانشین فنکشن نوں وار وار
	لاگو کر کے $0$ توں حاصل ہوندا اے۔
\end{enumerate}
پہلیاں دو شرطاں نوں پہلے درجے دی منطق نال سنبھالنا سوکھا اے (اُتّے
ویکھو)۔ پر صرف پہلے درجے دی منطق نال اسی
\olref{repeatedapplication} نوں نہیں سنبھال سکدے۔ ڈیڈیکائنڈ دی وڈی
پیش رفت ایہ سی کہ شرط \olref{repeatedapplication} نوں سیٹ تھیوری دے
لفظاں وچ ایویں نواں بیان کیتا جاوے:
\begin{enumerate}
	\item[3$'$.] قدرتی عدد اوہ سب توں چھوٹا سیٹ نیں جیہڑا
	جانشین فنکشن \emph{تحت بند} اے: یعنی جے اسی سیٹ دے کسے وی
	!!{element} اُتے $s$ لاگو کریے، تاں سانوں سیٹ دا اک ہور
	!!{element} ملدا اے۔
\end{enumerate}
پر سانوں ایہ گل ہولی ہولی کھولنی پینی اے۔

\begin{defn}\ollabel{Closure}
	اپنے ہی حامل سیٹ اندر نقشہ کردے ہر فنکشن $f$ لئی، اوس حامل دا ذیلی سیٹ $X$،
	$f$-\emph{بند} {اُدوں تے صرف اُدوں} اے جدوں
	$(\forall x \in X)f(x) \in X$۔ ہُن اوس حامل دے ہر $o$ لئی تعریف کرو:
	$$\closureofunder{f}{o} = \bigcap\Setabs{X}{o \in X\text{ تے }X
\text{، $f$-بند اے}}$$
\end{defn}
% PNBINF-002 ماخذ دی درستی: تعریف وچ f نوں اپنے حامل سیٹ دا خود نقشہ، X نوں حامل دا ذیلی سیٹ تے o نوں اوس حامل دا رکن سمجھیا اے؛ ایس دائرے بناں اگلی سطر دا A بے تعارف تے حاصل قدراں دا گواہ لازماً بند نہیں۔

سو $\closureofunder{f}{o}$ اوہناں سارے $f$-بند سیٹاں دا اشتراک اے
جیہناں وچ $o$، !!a{element} اے۔ فیر بدیہی طور تے
$\closureofunder{f}{o}$ اوہ \emph{سب توں چھوٹا} $f$-بند سیٹ اے جیہدے
وچ $o$، !!a{element} اے۔ اگلا نتیجہ ایس بدیہی خیال نوں پوری درستی
دیندا اے:
\begin{lem}\ollabel{closureproperties}
	ہر ایسے فنکشن $f$ تے ہر $o \in A$ لئی، جتھے A اوس فنکشن دا حامل سیٹ اے:
	\begin{enumerate}
		\item\ollabel{closurehaselem} $o \in \closureofunder{f}{o}$؛ تے
		\item\ollabel{closureclosed} $\closureofunder{f}{o}$، $f$-بند اے؛ تے
		\item\ollabel{closuresmallest} جے $X$، $f$-بند اے تے $o \in
		X$، تاں $\closureofunder{f}{o} \subseteq X$
	\end{enumerate}
\end{lem}

\begin{proof}
نوٹ کرو کہ گھٹ توں گھٹ اک $f$-بند سیٹ اے جیہدے وچ $o$،
!!a{element} اے، یعنی $\ran{f}\cup \{o\}$۔ سو
$\closureofunder{f}{o}$، ایسے \emph{سارے} سیٹاں دا اشتراک، موجود اے۔ ہُن سانوں
\olref{closurehaselem}--\olref{closuresmallest} جانچنے نیں۔

\olref{closurehaselem} بارے: $o \in \closureofunder{f}{o}$، کیوں جو
ایہ اوہناں سیٹاں دا اشتراک اے جیہناں ساریاں وچ $o$، !!a{element} اے۔

\olref{closureclosed} بارے: فرض کرو $x \in \closureofunder{f}{o}$۔
سو جے $o \in X$ تے $X$، $f$-بند اے، تاں $x \in X$، تے ہُن
$f(x) \in X$، کیوں جو $X$، $f$-بند اے۔ سو
$f(x) \in \closureofunder{f}{o}$۔

\olref{closuresmallest} بارے: عام طور تے، جے $X
\in C$، تاں $\bigcap C \subseteq X$۔
\end{proof}

ایہنوں ورت کے اسی آکھ سکدے آں:

\begin{defn}
اک \emph{ڈیڈیکائنڈ الجبرا} اک سیٹ $A$ اے جیہدے نال فنکشن $f
\colon A \to A$ تے کوئی $o \in A$ ایویں دِتے گئے نیں کہ:
	\begin{enumerate}
		\item \ollabel{ded:proper} $o \notin \ran{f}$
		\item \ollabel{ded:injection} $f$ اک !!a{injection} اے
		\item \ollabel{ded:closure} $A = \closureofunder{f}{o}$
	\end{enumerate}
\end{defn}

کیوں جو $A = \closureofunder{f}{o}$، ساڈا پہلا نتیجہ دَسدا اے کہ
$A$ اوہ سب توں چھوٹا $f$-بند سیٹ اے جیہدے وچ $o$، !!a{element} اے۔
صاف اے کہ ڈیڈیکائنڈ الجبرا ڈیڈیکائنڈ لامتناہی اے؛ تعریف دیاں شرطاں
\olref{ded:proper} تے \olref{ded:injection} ویکھ لو۔ پر ہور دلچسپ حقیقت
ایہ اے کہ ہر ڈیڈیکائنڈ لامتناہی سیٹ نوں ڈیڈیکائنڈ الجبرا بنایا جا
سکدا اے۔

\begin{thm}\ollabel{thm:DedekindInfiniteAlgebra}
جے کوئی ڈیڈیکائنڈ لامتناہی سیٹ اے، تاں کوئی ڈیڈیکائنڈ الجبرا اے۔
\end{thm}

\begin{proof}
$D$ نوں ڈیڈیکائنڈ لامتناہی لو۔ سو اک انجیکشن $g \colon D \to D$ تے
اک رکن $o  \in D \setminus \ran{g}$ اے۔ ہُن $A =
\closureofunder{g}{o}$ لو؛ \olref{closureproperties} مطابق $A$ موجود اے
تے $o \in A$۔ $f = \funrestrictionto{g}{A}$ لو۔ اسی وکھاواں گے کہ
$A, f, o$ رل کے اک ڈیڈیکائنڈ الجبرا بناؤندے نیں۔

\olref{ded:proper} بارے: $o \notin \ran{g}$ تے $\ran{f}
\subseteq \ran{g}$، سو $o\notin \ran{f}$۔

\olref{ded:injection} بارے: $g$، $D$ اُتے اک انجیکشن اے؛ سو $f
\subseteq g$ لازماً اک انجیکشن اے۔

\olref{ded:closure} بارے: \olref{closureproperties} مطابق $A$،
$g$-بند اے؛ ایس توں ہور وی صاف اے کہ $A$، $f$-بند اے۔ سو
$\closureofunder{f}{o} \subseteq A$، \olref{closureproperties} مطابق۔
کیوں جو $\closureofunder{f}{o}$ وی $f$-بند اے تے
$f = \funrestrictionto{g}{A}$، ایس لئی $\closureofunder{f}{o}$،
$g$-بند اے۔ سو $A \subseteq \closureofunder{f}{o}$،
\olref{closureproperties} مطابق۔
\end{proof}

\end{document}
